% GENERATED FILE. Edit canonical lessons, then run npm run generate:books.
% canonical-source-hash: fnv1a64:515cfd0637268b18
% canonical-lessons: ML-C434-pathyapaddhati, ML-C434-hajar, ML-C434-hosrral, ML-C434-skolarsippu, ML-C434-pasakuka

\chapter{Campus Life}
\label{ch:ml-campus-life}
\hlchaptermodality{434}

\begin{chapteropening}
\textbf{By the end of this chapter:} \emph{I can say a curriculum, attendance, a hostel, a scholarship and to pass an exam.}

Complete the last lesson of chapter 434: I can say a curriculum, attendance, a hostel, a scholarship and to pass an exam.
\end{chapteropening}

\section[\texorpdfstring{\ml{പാഠ്യപദ്ധതി} (pāṭhyapaddhati)}{പാഠ്യപദ്ധതി (pāṭhyapaddhati)}]{\ml{പാഠ്യപദ്ധതി} (pāṭhyapaddhati) — a curriculum}

\label{lesson:ML-C434-pathyapaddhati}



\begin{quote}
Before the new one: say the Malayalam for a semester, then the Malayalam for a syllabus.
\end{quote}

\subsection*{You'll want to know: \ml{പാഠ്യപദ്ധതി}}
\textbf{\ml{പാഠ്യപദ്ധതി}} — \emph{pāṭhyapaddhati} — \textquotedblleft{}a curriculum\textquotedblright{}.

\begin{quote}
\ml{സർക്കാർ} \ml{പുതിയ} \ml{പാഠ്യപദ്ധതി} \ml{കൊണ്ടുവന്നു}.
\end{quote}

\begin{quote}
\emph{sarkkār putiya pāṭhyapaddhati koṇṭuvannu.}
\end{quote}

\begin{quote}
The government brought in a new curriculum.
\end{quote}

\ml{പദ്ധതി} alone is a plan or a scheme.

\begin{cousinweb}
\textbf{\ml{പാഠ്യപദ്ധതി}} (\emph{pāṭhyapaddhati}) is a Malayalam compound built from Sanskrit \emph{pāṭhya}, \textquotedblleft{}to be recited, to be taught\textquotedblright{} (Monier-Williams), and \emph{paddhati}, \textquotedblleft{}a way, a path\textquotedblright{} (Monier-Williams): the path of what is taught.
\end{cousinweb}

\begin{grammarlens}[title={What you've built}]
One more word for giving reasons and opinions.
\end{grammarlens}

\subsection*{Your turn}
\begin{itemize}
\raggedright
  \item \emph{Say it:} \emph{pāṭhyapaddhati}
  \item \emph{Say it:} \emph{pāṭhyapaddhati}, once more
  \item \emph{Say it:} \emph{sarkkār putiya pāṭhyapaddhati koṇṭuvannu.} — The government brought in a new curriculum.
  \item \emph{Recall:} say the Malayalam for a semester, then the Malayalam for a syllabus, then say \emph{pāṭhyapaddhati} again
\end{itemize}

\begin{tcolorbox}[breakable,colback=teal!4,colframe=teal!35!black,title={Before you move on}]
What does \ml{പാഠ്യപദ്ധതി} mean? (\textquotedblleft{}A curriculum\textquotedblright{}.) Say it once more.
\end{tcolorbox}
\section[\texorpdfstring{\ml{ഹാജർ} (hājar)}{ഹാജർ (hājar)}]{\ml{ഹാജർ} (hājar) — attendance}

\label{lesson:ML-C434-hajar}



\begin{quote}
Before the new one: say the Malayalam for a syllabus, then the Malayalam for a curriculum.
\end{quote}

\subsection*{You'll want to know: \ml{ഹാജർ}}
\textbf{\ml{ഹാജർ}} — \emph{hājar} — \textquotedblleft{}attendance\textquotedblright{}.

\begin{quote}
\ml{ക്ലാസിൽ} \ml{ഇന്ന്} \ml{ഹാജർ} \ml{കുറവായിരുന്നു}.
\end{quote}

\begin{quote}
\emph{klāsil innŭ hājar kuṟavāyirunnu.}
\end{quote}

\begin{quote}
Attendance in class was low today.
\end{quote}

\ml{ഹാജരാകുക}, to be present, is built on it.

\begin{cousinweb}
\textbf{\ml{ഹാജർ}} (\emph{hājar}) came from Arabic \emph{ḥāḍir}, \textquotedblleft{}present\textquotedblright{}, and was borrowed from Hindustani \emph{hāzir}; it is the noun inside \ml{ഹാജരാകുക}.
\end{cousinweb}

\begin{grammarlens}[title={What you've built}]
One more word for giving reasons and opinions.
\end{grammarlens}

\subsection*{Your turn}
\begin{itemize}
\raggedright
  \item \emph{Say it:} \emph{hājar}
  \item \emph{Say it:} \emph{hājar}, once more
  \item \emph{Say it:} \emph{klāsil innŭ hājar kuṟavāyirunnu.} — Attendance in class was low today.
  \item \emph{Recall:} say the Malayalam for a syllabus, then the Malayalam for a curriculum, then say \emph{hājar} again
\end{itemize}

\begin{tcolorbox}[breakable,colback=teal!4,colframe=teal!35!black,title={Before you move on}]
What does \ml{ഹാജർ} mean? (\textquotedblleft{}Attendance\textquotedblright{}.) Say it once more.
\end{tcolorbox}
\section[\texorpdfstring{\ml{ഹോസ്റ്റൽ} (hōsṟṟal)}{ഹോസ്റ്റൽ (hōsṟṟal)}]{\ml{ഹോസ്റ്റൽ} (hōsṟṟal) — a hostel}

\label{lesson:ML-C434-hosrral}



\begin{quote}
Before the new one: say the Malayalam for a curriculum, then the Malayalam for attendance.
\end{quote}

\subsection*{You'll want to know: \ml{ഹോസ്റ്റൽ}}
\textbf{\ml{ഹോസ്റ്റൽ}} — \emph{hōsṟṟal} — \textquotedblleft{}a hostel\textquotedblright{}.

\begin{quote}
\ml{അവൾ} \ml{കോളേജ്} \ml{ഹോസ്റ്റലിൽ} \ml{താമസിക്കുന്നു}.
\end{quote}

\begin{quote}
\emph{avaḷ kōḷējŭ hōsṟṟalil tāmasikkunnu.}
\end{quote}

\begin{quote}
She stays in the college hostel.
\end{quote}

Plural: \ml{ഹോസ്റ്റലുകൾ}.

\begin{cousinweb}
\textbf{\ml{ഹോസ്റ്റൽ}} (\emph{hōsṟṟal}) was borrowed from English \emph{hostel}.
\end{cousinweb}

\begin{grammarlens}[title={What you've built}]
One more word for giving reasons and opinions.
\end{grammarlens}

\subsection*{Your turn}
\begin{itemize}
\raggedright
  \item \emph{Say it:} \emph{hōsṟṟal}
  \item \emph{Say it:} \emph{hōsṟṟal}, once more
  \item \emph{Say it:} \emph{avaḷ kōḷējŭ hōsṟṟalil tāmasikkunnu.} — She stays in the college hostel.
  \item \emph{Recall:} say the Malayalam for a curriculum, then the Malayalam for attendance, then say \emph{hōsṟṟal} again
\end{itemize}

\begin{tcolorbox}[breakable,colback=teal!4,colframe=teal!35!black,title={Before you move on}]
What does \ml{ഹോസ്റ്റൽ} mean? (\textquotedblleft{}A hostel\textquotedblright{}.) Say it once more.
\end{tcolorbox}
\section[\texorpdfstring{\ml{സ്കോളർഷിപ്പ്} (skōḷarṣippŭ)}{സ്കോളർഷിപ്പ് (skōḷarṣippŭ)}]{\ml{സ്കോളർഷിപ്പ്} (skōḷarṣippŭ) — a scholarship}

\label{lesson:ML-C434-skolarsippu}



\begin{quote}
Before the new one: say the Malayalam for attendance, then the Malayalam for a hostel.
\end{quote}

\subsection*{You'll want to know: \ml{സ്കോളർഷിപ്പ്}}
\textbf{\ml{സ്കോളർഷിപ്പ്}} — \emph{skōḷarṣippŭ} — \textquotedblleft{}a scholarship\textquotedblright{}.

\begin{quote}
\ml{നല്ല} \ml{മാർക്ക്} \ml{കിട്ടിയതിനാൽ} \ml{അവന്} \ml{സ്കോളർഷിപ്പ്} \ml{കിട്ടി}.
\end{quote}

\begin{quote}
\emph{nalla mārkkŭ kiṭṭiyatināl avanŭ skōḷarṣippŭ kiṭṭi.}
\end{quote}

\begin{quote}
Because he got good marks, he got a scholarship.
\end{quote}

A scholarship is \ml{കിട്ടുക}, received.

\begin{cousinweb}
\textbf{\ml{സ്കോളർഷിപ്പ്}} (\emph{skōḷarṣippŭ}) was borrowed from English \emph{scholarship}.
\end{cousinweb}

\begin{grammarlens}[title={What you've built}]
One more word for giving reasons and opinions.
\end{grammarlens}

\subsection*{Your turn}
\begin{itemize}
\raggedright
  \item \emph{Say it:} \emph{skōḷarṣippŭ}
  \item \emph{Say it:} \emph{skōḷarṣippŭ}, once more
  \item \emph{Say it:} \emph{nalla mārkkŭ kiṭṭiyatināl avanŭ skōḷarṣippŭ kiṭṭi.} — Because he got good marks, he got a scholarship.
  \item \emph{Recall:} say the Malayalam for attendance, then the Malayalam for a hostel, then say \emph{skōḷarṣippŭ} again
\end{itemize}

\begin{tcolorbox}[breakable,colback=teal!4,colframe=teal!35!black,title={Before you move on}]
What does \ml{സ്കോളർഷിപ്പ്} mean? (\textquotedblleft{}A scholarship\textquotedblright{}.) Say it once more.
\end{tcolorbox}
\section[\texorpdfstring{\ml{പാസാകുക} (pāsākuka)}{പാസാകുക (pāsākuka)}]{\ml{പാസാകുക} (pāsākuka) — to pass an exam}

\label{lesson:ML-C434-pasakuka}



\begin{quote}
Before the new one: say the Malayalam for a hostel, then the Malayalam for a scholarship.
\end{quote}

\subsection*{You'll want to know: \ml{പാസാകുക}}
\textbf{\ml{പാസാകുക}} — \emph{pāsākuka} — \textquotedblleft{}to pass an exam\textquotedblright{}.

\begin{quote}
\ml{അവൾ} \ml{എല്ലാ} \ml{വിഷയങ്ങളിലും} \ml{പാസായി}.
\end{quote}

\begin{quote}
\emph{avaḷ ellā viṣayaṅṅaḷiluṁ pāsāyi.}
\end{quote}

\begin{quote}
She passed in all subjects.
\end{quote}

Past: \ml{പാസായി}, \emph{pāsāyi}.

Its past is \ml{പാസായി}; the causative \ml{പാസാക്കുക} was taught a few chapters earlier.

\begin{cousinweb}
\textbf{\ml{പാസാകുക}} (\emph{pāsākuka}) is a Malayalam compound: the first part was borrowed from English \emph{pass}, and \ml{ആകുക}, \textquotedblleft{}to become\textquotedblright{}, is native Dravidian (DEDR 333).
\end{cousinweb}

\begin{grammarlens}[title={What you've built}]
One more word for giving reasons and opinions.
\end{grammarlens}

\subsection*{Your turn}
\begin{itemize}
\raggedright
  \item \emph{Say it:} \emph{pāsākuka}
  \item \emph{Say it:} \emph{pāsākuka}, once more
  \item \emph{Say it:} \emph{avaḷ ellā viṣayaṅṅaḷiluṁ pāsāyi.} — She passed in all subjects.
  \item \emph{Recall:} say the Malayalam for a hostel, then the Malayalam for a scholarship, then say \emph{pāsākuka} again
\end{itemize}

\begin{tcolorbox}[breakable,colback=teal!4,colframe=teal!35!black,title={Before you move on}]
What does \ml{പാസാകുക} mean? (\textquotedblleft{}To pass an exam\textquotedblright{}.) Say it once more.
\end{tcolorbox}
